%% ============================================================
%%  Piezas comunes a los dos decks: Best Model Overview y glosario
%% ============================================================

% ---- Best Model Overview (diapositiva obligatoria, primera de cada deck) ----
\newcommand{\diapositivaoverview}{%
\begin{frame}{Best Model Overview}
  {Ingreso predicho por persona + LightGBM \;|\; Equipo 06, MECA 4107}
  \vspace{-0.1cm}
  \begin{columns}[T]
    \column{0.37\textwidth}\centering \keyfig{\kagglepublico}{F1 público en Kaggle}
    \column{0.30\textwidth}\centering \keyfig{0,782}{F1 en holdout, IC 95\,\% [0,774; 0,789]}
    \column{0.30\textwidth}\centering \keyfig{0,776}{F1 en CV repetida (15 evaluaciones)}
  \end{columns}
  \vspace{0.15cm}
  \footnotesize
  \begin{columns}[T]
    \column{0.49\textwidth}
      \textbf{Especificación}
      \begin{itemize}
        \item \textbf{Variable clave:} el ingreso de cada persona, predicho con lo que el test
          sí trae y sumado al hogar (1.324 variables en total).
        \item \textbf{Algoritmo:} LightGBM (76\,\%) mezclado con dos redes (19\,\%) y un
          rediseño desde cero (5\,\%); pesos elegidos fuera de muestra.
        \item \textbf{Regla:} pobre si $P(\text{pobre}) \geq 0{,}36$.
      \end{itemize}
    \column{0.49\textwidth}
      \textbf{Mensajes}
      \begin{enumerate}
        \item Las variables aportaron \textbf{9 veces más} F1 que los algoritmos.
        \item El techo (${\approx}\,0{,}78$) viene sobre todo de la información que falta.
        \item Excluye al 18\,\% de los pobres, casi todos cerca de la línea.
      \end{enumerate}
  \end{columns}
  \fuente{F1: media armónica de precisión y recall de la clase pobre. Holdout: 20\,\% del train nunca usado para decidir. Fuente: views/tablas/nivel7/etapa18\_resumen.csv.}
  \note{
    \textbf{Qué decir (1 min).} Nuestro mejor modelo clasifica hogares pobres con F1 de
    0,782 en datos que nunca usamos para decidir. La pieza clave no es el algoritmo sino
    una variable: predecimos cuánto gana cada persona con lo que el test sí trae
    (ocupación, horas, cotización) y lo sumamos al hogar. Tres mensajes: las variables
    pesan 9 veces más que el algoritmo, hay un techo que viene de los datos y los errores
    se concentran cerca de la línea.\\[3pt]
    \textbf{Pregunta difícil: ¿por qué una mezcla y no un solo modelo?} La mezcla gana
    +0,0007 en CV repetida frente a LightGBM solo. El IC ingenuo [+0,0002; +0,0012] excluye
    cero; el corregido por dependencia entre folds no. Es una mejora marginal: en la
    práctica LightGBM solo rinde igual.
  }
\end{frame}}

% ---- Glosario (apéndice de ambos decks, dos diapositivas) ----
\newcommand{\filaglosario}[2]{\textbf{#1} & #2 \\}
\newcommand{\diapositivaglosario}{%
\begin{frame}{Glosario (1/2): métricas y errores}
  \footnotesize\renewcommand{\arraystretch}{1.25}
  \begin{tabular}{@{}p{0.2\textwidth}p{0.76\textwidth}@{}}
    \toprule
    \filaglosario{Precisión}{De los hogares que el programa atiende, qué parte es pobre. $1-$precisión $=$ \textbf{filtración}: no pobres que reciben el beneficio.}
    \filaglosario{Recall}{De los hogares pobres, qué parte atiende el programa. $1-$recall $=$ \textbf{subcobertura}: pobres excluidos.}
    \filaglosario{F1}{Media armónica de precisión y recall: solo es alta si ambas lo son. No premia a los no pobres bien clasificados.}
    \filaglosario{Umbral}{Probabilidad a partir de la cual un hogar se clasifica como pobre. Bajarlo cubre más pobres a cambio de más filtración.}
    \filaglosario{IC 95\,\%}{Rango de valores compatibles con los datos. Si el IC de una diferencia incluye cero, no podemos afirmar que haya mejora.}
    \filaglosario{Efecto mínimo detectable}{La mejora más pequeña que el diseño detectaría con probabilidad 0,8.}
    \bottomrule
  \end{tabular}
  \note{Glosario de apoyo; se muestra solo si lo preguntan.}
\end{frame}
\begin{frame}{Glosario (2/2): validación y variables}
  \footnotesize\renewcommand{\arraystretch}{1.25}
  \begin{tabular}{@{}p{0.2\textwidth}p{0.76\textwidth}@{}}
    \toprule
    \filaglosario{Holdout}{20\,\% del train (32.993 hogares) apartado al inicio; solo se usa para medir, nunca para decidir.}
    \filaglosario{Validación cruzada}{El 80\,\% restante se parte en 5 bloques (folds): se entrena con 4 y se evalúa en el quinto, rotando.}
    \filaglosario{Fuera de fold (OOF)}{La predicción de cada hogar hecha por un modelo que no lo vio. Con ella se eligen umbral, pesos e hiperparámetros.}
    \filaglosario{CV repetida}{La validación cruzada repetida con 3 particiones nuevas: 15 evaluaciones para medir la variabilidad.}
    \filaglosario{Auxiliar de ingreso}{Modelo que predice el ingreso de cada persona con variables que el test sí trae; sumado al hogar, entra como variable.}
    \filaglosario{OOF anidada}{El auxiliar que alimenta a un hogar no vio ese hogar ni su fold de validación: así no se filtra la respuesta.}
    \bottomrule
  \end{tabular}
  \note{Glosario de apoyo; se muestra solo si lo preguntan.}
\end{frame}}
